\section{Positive Control: Is Certified Correctness Recoverable at All?}
\label{sec:control}

A negative result for shared operators says little unless some intervention
can move the certificate in the same regime. We therefore report a separate, exploratory, supervised
\FarlaShots-shot class-specific pilot, originally developed as a candidate class-specific method and
used here only as a positive control, on fresh reserved items of the same two
backbones, datasets and noise level, with \FarlaCifarConfirmItems{} CIFAR-100
and \FarlaEurosatConfirmItems{} EuroSAT confirmation items per cell and the same
certification budget. Using \FarlaShots{} labeled images per class, it trained
class-specific low-rank tangent adjustments of the prototype bank in the spirit
of parameter-efficient adaptation~\cite{hu2022lora} under a noisy-and-clean
cross-entropy objective. Seven variants were evaluated (Table~\ref{tab:farla}):
a supervised \emph{shared} translation (one common correction, the
comparator); a \emph{low-rank} adapter with the cross-entropy objective only; a
\emph{tail-only} variant that adds a loss on the vote fraction near the
certification threshold; the \emph{full} objective at rank 8, at rank 4, and in
a direct (non-tangent) parametrization; and a \emph{pseudo-label} variant of the
full objective trained on the model's own clean predictions instead of labels.
The pilot was registered with its own configuration and
custody package, but its planned smoke run was skipped, so we call it
exploratory.

% Generated by scripts/generate_satml2027_paper_assets.py; do not edit by hand.
% Positive-control pilot candidates and planned contrasts.
% Sources (SHA-256):
%   FARLA/downloaded_results/EXP-20260917-020-FARLA-FULL/analysis/report.json  dd123e618475ebb498583fb569e26040f45cf051b6d8338224d61af37a2777c4
%   FARLA/downloaded_results/EXP-20260917-020-FARLA-FULL/config_snapshot.json  57b79923fd333c01fb128e7dc88f81f9aefce7e6b3aafff4d373d5c5d5511620
%   FARLA/downloaded_results/EXP-20260917-020-FARLA-FULL/splits/manifest.json  cdda5225e1220b207d084a57a23c6771e396c0ba5dda57d3f7223710097d960d
\begin{table*}[t]
\centering
\scriptsize
\setlength{\tabcolsep}{4pt}
\caption{Positive-control pilot: macro standard CA, anchored CA and certified-wrong rate at radius 0.25 and macro clean accuracy over 4 cells (1,000 CIFAR-100 and 100 EuroSAT confirmation items per cell, 20 supervised shots per class), with all five planned paired contrasts and their class-stratified percentile bootstrap 95\% intervals (10,000 replicates). Clean gate: macro drop at most 1 point and every cell drop at most 2 points.}
\label{tab:farla}
\begin{tabular}{lrrrrc}
\toprule
Candidate & Std CA & Anch.\ CA & Clean & Cert.-wrong & Clean gate \\
\midrule
No correction & 11.60 & 8.43 & 60.88 & 34.18 & pass \\
Shared translation (noisy CE) & 11.75 & 8.68 & 61.28 & 33.43 & pass \\
Low-rank r=8 (noisy CE) & 27.10 & 23.35 & 72.50 & 13.58 & pass \\
Tail-only r=8 & 36.50 & 28.75 & 70.80 & 9.05 & pass \\
Full objective r=4 & 33.35 & 24.60 & 62.15 & 11.78 & pass \\
Full objective r=8 & 35.05 & 29.55 & 68.58 & 10.43 & pass \\
Full objective, direct & 34.60 & 31.20 & 66.23 & 11.43 & fail \\
Full objective r=8, pseudo-labels (no ground truth) & 27.10 & 21.15 & 46.98 & 17.75 & fail \\
\bottomrule
\end{tabular}
\par\vspace{4pt}
\begin{tabular}{lrrr}
\toprule
Planned contrast (Std CA) & $\Delta$ & 95\% lower & 95\% upper \\
\midrule
Full objective r=8 vs.\ no correction & $+23.45$ & $+20.25$ & $+26.63$ \\
Full objective r=8 vs.\ tail-only r=8 & $-1.45$ & $-4.05$ & $+1.13$ \\
Full objective r=8 vs.\ full objective r=4 & $+1.70$ & $-0.65$ & $+4.00$ \\
Full objective r=8 vs.\ low-rank r=8 (noisy CE) & $+7.95$ & $+4.73$ & $+11.15$ \\
Low-rank r=8 (noisy CE) vs.\ shared translation (noisy CE) & $+15.35$ & $+12.70$ & $+18.03$ \\
\bottomrule
\end{tabular}
\end{table*}


Table~\ref{tab:farla} shows that class-specific adaptation recovers a large
share of certified correctness in a regime where shared translations did not:
standard certified accuracy at $r=0.25$ rises from \FarlaIdentityStdCA\% to
\FarlaFullStdCA\% for the full rank-8 variant, a paired macro gain of
$\FarlaVsNoCorrDelta$ points with 95\% interval
$[\FarlaVsNoCorrLower,\FarlaVsNoCorrUpper]$, while the supervised shared
translation reaches only \FarlaSharedStdCA\%. The certified-wrong rate falls
from \FarlaIdentityCW\% to \FarlaFullCW\%. Four qualifications bound the
reading. The tail-only ablation reaches \FarlaTailStdCA\%, above the full variant's
point estimate, with a certified-wrong rate of \FarlaTailCW\%, and the full
variant's contrast against it is
$\FarlaVsTailDelta$ points with interval
$[\FarlaVsTailLower,\FarlaVsTailUpper]$; that interval does not establish a
difference, so the additional certificate-aware terms are not shown to be
necessary. Clean accuracy also rises, from
\FarlaIdentityClean\% to \FarlaFullClean\% (and to \FarlaTailClean\% for the
tail-only variant), so these results do not isolate a benefit specific to
certificate targeting from ordinary accuracy improvement under supervised
adaptation. The pseudo-label variant, the only one that uses no
ground truth, reaches \FarlaPseudoStdCA\% standard certified accuracy
($\FarlaPseudoVsNoCorrDelta$ points,
$[\FarlaPseudoVsNoCorrLower,\FarlaPseudoVsNoCorrUpper]$) but lowers clean
accuracy to \FarlaPseudoClean\%, a drop of \FarlaPseudoCleanDrop{} points that
fails the clean gate: it buys certified stability at the cost of clean
correctness, which is the trade this paper argues must be reported explicitly.
Finally, the pilot itself was not compared with noise-aware prompt
learning~\cite{hussein2024promptsmooth}, the natural baseline for any text-side
adaptation; the follow-up adds a few-shot prompt control inspired by it
(Section~\ref{sec:followup}). The control therefore establishes a scope boundary: the certificate
in this regime is movable by class-specific supervised adaptation, and the
result of Section~\ref{sec:results} concerns shared operators only.
